\section{Method}
\label{sec:method}

\subsection{Task and encoding}

We train on $c = (a+b) \bmod P$ with $P = 113$, over the full grid of $P^2 =
12\,769$ ordered pairs. The input is the concatenation of two one-hot vectors,
$x = [\,\mathrm{onehot}(a) \,\|\, \mathrm{onehot}(b)\,] \in \R^{2P}$, and the
output is a distribution over the $P$ possible values of $c$. No embedding
layer is used, following \citet{bussmann2023modadd}.

This encoding is chosen deliberately rather than for convenience. Because the
first $P$ coordinates of $x$ depend only on $a$ and the next $P$ only on $b$,
any weight matrix acting on the input decomposes into a block that is a
function of $a$ and a block that is a function of $b$. Its discrete Fourier
transform along either block is therefore directly interpretable, with no
probing, no auxiliary model, and no approximation. The same applies to router
keys, which under one-hot inputs reduce to learned lookup tables indexed by
the operand value.

A second dataset is used as a positive control in
Section~\ref{sec:experiments}: examples are drawn from two subtasks,
$(a+b) \bmod P$ and $(a \cdot b) \bmod P$, with a one-hot task token appended
to the input. Here a genuine two-way partition of the input distribution
exists, so expert specialisation by subtask is possible in principle and can
be detected. Its role is to distinguish a negative specialisation result on
the single task from an artefact of the measurement.

\subsection{Architectures}

Every model consists of a single hidden layer mapping $\R^{2P} \to \R^{P}$.
The baseline is a dense two-layer MLP with $100$ hidden units and ReLU. The
six Mixture-of-Experts variants share one product-key router, defined in
Equation~\eqref{eq:product-key}, with $n$ keys per side, $N = n^2$ composed
experts, four heads, and top-$k$ selection with $k = 4$. They differ only in
how the expert tensor is represented: unfactorised, MONET-HD, MONET-VD, CP
(rank $128$), Tucker (rank $32$) and Tensor-Train (rank $16$). Expert
dimension is $m = 8$ throughout.

We sweep $n \in \{4, 8, 16, 32, 64\}$, that is $N$ from $16$ to $4096$. The
unfactorised layer is capped at $n \le 32$, since its parameter count grows
linearly in $N$ and at $n = 64$ it would hold $11.7$M parameters against
$239$K for MONET at the same expert count; beyond that point the comparison
measures capacity rather than structure.

Two deviations from \citet{park2025monet} are worth stating. First, MONET
estimates routing quantiles with BatchNorm to avoid a hardware-unfriendly
top-$k$; at $N \le 4096$ an exact top-$k$ is inexpensive, so we use it, and
retain BatchNorm as an option. Second, the auxiliary losses below are computed
on the full softmax rather than on the top-$k$-masked weights, which contain
exact zeros.

\subsection{Training protocol}

All models are trained full-batch with AdamW, learning rate $10^{-3}$,
$\beta = (0.9, 0.98)$, weight decay $1.0$, gradient clipping at norm $1.0$,
for $30\,000$ steps. The training split is $80\%$ of the grid, chosen once per
dataset and held fixed across all runs and seeds; Section~\ref{sec:protocol}
reports the sweep that led to this choice.

Following MONET we add a uniformity and an ambiguity term to the
cross-entropy,
\begin{equation}
  \mathcal{L}
  = \mathcal{L}_{\mathrm{CE}}
  + \lambda\big(\mathcal{L}_{\mathrm{unif}} + \mathcal{L}_{\mathrm{amb}}\big),
  \qquad \lambda = 10^{-3},
\end{equation}
where $\mathcal{L}_{\mathrm{unif}}$ penalises routing distributions that
deviate from uniform and $\mathcal{L}_{\mathrm{amb}}$ penalises routing
distributions that are not peaked, the two acting in opposition.

Weight decay is applied to weight matrices and, importantly, to the router
keys; biases and normalisation parameters are exempt. This is not the usual
convention, and Section~\ref{sec:ablations} shows it is necessary: under
one-hot inputs the router keys form a lookup table indexed by the operand, so
leaving them unregularised hands the model an unpenalised memorisation
mechanism.

The seed controls weight initialisation only. The train/test partition is
identical across seeds, so the reported spread measures variation due to
initialisation at fixed data, not variation due to resampling. We use three
seeds for every configuration that carries a claim.

\subsection{Metrics}

\paragraph{Grokking step.} The first optimisation step at which test accuracy
reaches $95\%$, evaluated every $100$ steps. Runs that never reach it within
the budget are reported as failures rather than assigned a censored value.

\paragraph{Spectral purity.} Let $f \in \R^{P}$ be the response of a unit as a
function of one operand. We remove its mean, take the discrete Fourier
transform, and normalise the power spectrum over frequencies $k = 1, \ldots,
(P-1)/2$ so that it sums to one. Purity is the share of the strongest
component,
\begin{equation}
  \mathrm{purity}(f) = \max_k \, \hat{p}_k .
\end{equation}
A value of $1$ means the unit is a pure single-frequency detector, which is
what monosemanticity means on this task. The DC component is dropped because a
constant offset carries no information about which residue class a unit
responds to.

When averaging across units we weight each unit by the $L_2$ norm of its
mean-removed response. Unused units still produce a noise spectrum, and an
unweighted mean is dragged down as soon as a layer leaves capacity idle;
weighting measures the purity of the units that actually carry signal. As a
reference point, white-noise units of the same dimension score $0.082$, and
all purity values below should be read against that floor.

\paragraph{Effective number of frequencies.} Purity describes single units;
the layer as a whole is described by pooling the normalised spectra of all its
units into $\bar{p}$ and taking the participation ratio
\begin{equation}
  n_{\mathrm{eff}} = \Big(\textstyle\sum_k \bar{p}_k^2\Big)^{-1}.
\end{equation}
This counts how many frequencies the layer actually uses, without a threshold:
a single pure tone gives $1$, a flat spectrum gives $56$. We prefer it to
counting frequencies above a fixed share of the mass, which is
threshold-sensitive and degenerates for layers that spread power evenly.

\paragraph{Router--basis spectral overlap.} Both the shared basis and the
router keys yield pooled spectra, $\bar{p}^{\mathrm{basis}}$ and
$\bar{p}^{\mathrm{router}}$. Their overlap is the total variation-style
intersection
\begin{equation}
  \mathrm{overlap} = \sum_k \min\!\big(\bar{p}^{\mathrm{basis}}_k,\,
  \bar{p}^{\mathrm{router}}_k\big) \in [0,1],
\end{equation}
which is $1$ when the two components operate on identical frequencies and $0$
when they are disjoint. Independent random spectra give roughly $0.7$, so
values well below that indicate active separation rather than chance.

\subsection{Reproducibility}

All analysis is performed on NumPy from archived weights; the router is
re-derived analytically rather than by running the model, which is exact under
one-hot inputs. The rearranged forward passes, which compute each factorised
layer without materialising the $N$-expert tensor, are verified against the
naive per-expert sum to machine precision for all six MoE variants; the test
is included in the repository. Code and aggregated run summaries are available at
\url{https://github.com/TheLastofDeadIns/Scaling-efficient-monosemantic-experts-in-transformers-with-tensor-decompositions}.